\chapter{Event Organizer Rules}

\section{Venue}

\begin{enumerate}[leftmargin=1.5em,labelsep=*]
	\item All Track Racing Disciplines, as well as Track Coasting and Track Gliding, must be held on a 
  standard athletics competition track. Exceptions are the 50m One Foot and the 30m Wheel Walk, 
  which can also be held at an indoor athletics venue. The IUF Slalom must be held on a surface 
  that has the same characteristics as a standard athletics competition track and therefore 
  complies with paragraph 2.1.
	If the requirements cannot be fulfilled, the disciplines cannot be offered as official 
  IUF Track Disciplines and this has to be announced accordingly before the competition.\\

	\textbf{Note:} If the track is outdoors, plans must be made to deal with inclement weather.
	Using an indoor track can eliminate this problem.
	The track must be available for enough days to allow for inclement weather.
	
	\item Every track and field competition venue for athletics approved by a national athletics 
  organization or higher is approved for official unicycle competitions.\\
	The essential criteria that such a venue must comply with are the following:
	\begin{enumerate}[label*=\arabic*,leftmargin=1.5em,labelsep=*]
		\item Any firm, uniform surface that complies with the specifications for syntetic surfaces for 
    athletic competition venues is permitted.
		
		\item The nominal length of a standard competition track is 400 m, it must be not shorter than 
    400.00 m and not longer than 400.04 m. It must consist of two parallel straights and two bends 
    whose radii must be equal. The inside of the track must be bordered by a kerb of suitable material 
    that should be coloured white, with a height of 50 mm to 65 mm and a width of 50 mm to 250 mm.
		The kerb on the two straights may be omitted and a white line 50 mm wide substituted.
		
		\item The measurement must be taken 0.30 m outward from the kerb or, where no kerb exists on a bend, 
    0.20 m from the line marking the inside of the track.
		
		\item The distance of the race must be measured from the edge of the start line farther from the 
    finish to the edge of the finish line nearer to the start.
		
		\item In all races up to and including 400 m, each rider must have a separate lane, with a width of 
    1.22 m $\pm$ 0.01 m, including the lane line on the right, marked by white lines 50 mm in width.
		All lanes must be of the same nominal width. The inner lane must be measured as stated in 
    pragraph 2.3, but the remaining lanes must be measured 0.20 m from the outer edges of the lines.
	\end{enumerate}
	
	\item The track must be prepared with start and finish lines for the various racing events that are 
  unique to unicycle racing (such as 50, 30, 10 and 5 meter lines).
	
	\item For races that include at least one turn, the inside lane must be on the left in the 
  direction of the race. The individual lanes must be numbered, starting with the left-hand lane as No. 1.
	
	\item A public address system must be provided to announce upcoming events and race winners.
	Bullhorns are usually not adequate for the track environment.
\end{enumerate}

\section{Officials and Judges}

\begin{enumerate}[leftmargin=1.5em,labelsep=*]
	\item The host must designate the following officials for all track events:
	\begin{itemize}
		\item Track Director
		\item Referee
	\end{itemize}
	\begin{enumerate}[label*=\arabic*,leftmargin=1.5em,labelsep=*]
		\item For races the following additional officials and judges must be designated:
	\end{enumerate}
	\begin{itemize}
		\item Starter
		\item Timekeeper
		\item Finish Line Judge
		\item Lane Judge
	\end{itemize}
	\begin{enumerate}[label*=\arabic*,leftmargin=1.5em,labelsep=*]
		\item For technical disciplines the following additional judges must be designated:
	\end{enumerate}
	\begin{itemize}
		\item Technical Disciplin Judge
	\end{itemize}
\end{enumerate}

\section{Communication}

\begin{comment2016}
Some thoughts on additional required communication:
\begin{itemize}
\item any changes to the standard rules
\item false start method
\item cut-off time
\item age groups
\item whether the event qualifies for a world record
\item whether helmets are required
\end{itemize}
\end{comment2016}

\begin{comment2016}
There should be a rule about when and how results are posted, so that the protest period can begin.
\end{comment2016}

\begin{enumerate}[leftmargin=1.5em,labelsep=*]
	\item If a convention host advertises events for disciplines with the names of the ones detailed in this chapter, 
  they must use the rules provided here.
	If hosts desire to do variations on these rules, for example by offering other unicycle classes or wheel sizes, 
  the events must be labeled accordingly, i.e. ``100m Unlimited'' or ``Track Coasting; Modified''.
	The host can also offer events for additional disciplines.
	In both cases, the events for modified disciplines and the additional disciplines, they cannot be 
  considered official IUF disciplines.
	In cases such as this, hosts must remember to provide detailed rules for these events at the same time the 
  events are announced.\\
	\textbf{Note:} Examples of modified discipline events would be Unlimited races, where races can be run on 
  unicycles without any restrictions. An example of another wheel size category would be the 
  700c wheel category, where unicycle wheels must be greater than 618mm in diameter, 
  have a maximum bead seat diameter (BSD) of 622 mm, and there are no restrictions on crank length.
	
	\item A Host is allowed to make helmets and/or other safety equipment mandatory for 
  the competition or individual disciplines but it must be announced when registration is opened and 
  must appear as an extra point to check when the competitor registers.
\end{enumerate}

\section{Age Groups \label{subsec:track_racing-categories_age-groups}}

\begin{enumerate}[leftmargin=1.5em,labelsep=*]
%TODO Ausnahme für die Staffeln hinzufügen
	\item The following age groups are the minimum required by the IUF to be offered at the time of 
  registration for any Track discipline: 0-10 (20 Class), 0-13, 14-18, 19-29, 30-UP.
	
	\item Convention hosts are free to offer more age groups, and often do.\\
	\textbf{Example:} A full range of offered age groups might look like 0-8 (20 Class), 9- 10 (20 Class), 
  0-12, 13-14, 15-16, 17-18, 19-29, 30-39, 40-49, 50-59, 60-UP.
	
%TODO Ausnahme für die Staffeln hinzufügen
	\item All age groups must be offered as men and women age group.
\end{enumerate}

\begin{comment2016}
\section{Practice}

Does there need to be a rule to allow practice on the track?
\end{comment2016}

\section{Cross-Discipline Rankings}

\subsection{Track Combined Competition}

\begin{enumerate}[leftmargin=1.5em,labelsep=*]
	\item The best finishers combined from 100 m, 400 m, 800 m, One Foot, Wheel Walk and IUF Slalom will win this title.
	\item Placement points are assigned for a placement in each of the above disciplines, based upon 
  finals/final rankings.
	The following system is used: 1st place gets 8, 2nd place 5, 3rd place 3, 4th place 2, and 5th place 1 point.
	Highest total points score is the Track Combined Champion; one each for mens and womens.
	\item If there is a tie, the rider with the most first places wins. If this still results	in a tie, 
  the title goes to the better finisher in the 100 m race.
\end{enumerate}

\subsection{Event-specific Overall Ranking}

\begin{enumerate}[leftmargin=1.5em,labelsep=*]
	\item It is up to the organizers to determine from which disciplines placement points will be awarded for the 
  Event-specific Overall Ranking.\\
	\textbf{Note:} The winner of this Event-specific Overall Ranking could be called for example 
  the "Overall Winner of XY Cup".
	\item Placement points are assigned in the individual disciplines.
    Usually the following system is used: 1st place gets 8, 2nd place 5, 3rd place 3, 4th place 2, 
    and 5th place 1 point.
    Highest total points score is the winner of the Event-specific Overall Ranking.
    Men and women should always be separated.
    There are two different ways of determining an overall ranking:
	
	\begin{enumerate}[label*=\arabic*,leftmargin=1.5em,labelsep=*]
		\item As an age group combined overall ranking, points will be awarded for placements in the 
    finals/final rankings, no points will be awarded in the age group rankings.
		
		\item As an age group separated overall ranking, points will be awarded for placements in the 
    corresponding age group, no points will be awarded in the final races/final rankings.
	\end{enumerate}
	
	\item If there is a tie, the rider with the most first places wins.
    If this still results in a tie, the title goes to the better finisher in the 100 m race or, 
    if this race is not part of the disciplines, it should go to the better finisher of the next 
    longer standard track race (according to 2B.2.2 - 2B.2.4).
    If there is no standard track race longer than 100 m, the organizers can use another predefined 
    discipline as tiebreaker.
	\item The choice of disciplines, the system to assign placement points and the tiebreaker discipline 
    should be published at the same time as the registration form or earlier.
\end{enumerate}

\section{Race Configuration}

\begin{enumerate}[leftmargin=1.5em,labelsep=*]
	\item Racing competition is held in two separate divisions: Men and Women. An exception are the relays, 
  which can be mixed men/women according to 2B.2.7, 7.

	\item There will be no mixing of age groups, or men and women (except for the relays), 
  in heats except with permission from the Racing Referee.

	\item Track events must have both a preliminary and final round.
\end{enumerate}

\section{Heat and Lane Assignments}

The following heat and lane assignments must be used for Unicon and international competitions.
Also for other competitions it is recommended to do the assignments accordingly.
The rule is applied for each age group independently.
\begin{enumerate}[leftmargin=1.5em,labelsep=*]
	\item The riders with the fastest seed times will be placed in the last heat, 
   the next riders in the second last heat, etc.\ until all riders are distributed over the heats.
	\item In age group races the distribution is done according to the seed times.
    Riders for whom no seed times are given will be placed without time behind the rider with the slowest seed time.
    The order in which riders with the same time are seeded will be decided by lot.
	\item In final races the distribution is done according to the times achieved in the age group races.
	The order in which riders with the same time will be seeded must be decided by lot.
	\item The lane assignment in lane-bound races is carried out according to the subsection below.
	\item In each heat at least three riders should be seated if possible; 
    however, this number can be undercut due to cancellations.
\end{enumerate}

\subsection{Lane assignments in lane-bound races}

\begin{enumerate}[leftmargin=1.5em,labelsep=*]
\item In races up to and including 100m the lanes are to be distributed as follows in each heat:
	\begin{itemize}
		\item[$-$] If the number of lanes is odd, the rider with the fastest seed time in the race will be placed on the middle lane.
		The rider with the next fastest seed time will be placed on the lane to the right of the middle lane (number of the middle lane +1) and all other riders will be placed alternately to the left and right of the middle lane according to their seed times.
		\item[$-$] If the number of lanes is even, the rider with the fastest seed time will be placed on the lane with a number value of half the total number of lanes.
		The rider with the next fastest seed time will be placed to the right of this lane (half lane number +1) and all other riders will be placed alternately to the left and right according to their seed times.
	\end{itemize}
	\item In races from 200m the lanes are to be distributed as follows in each heat:
	\begin{itemize}
		\item[$-$] The rider with the fastest seed time in the race will be placed on lane 1, the rider with the next fastest seed time will be placed on lane 2 and all other riders will be placed one lane higher according to their seed times.
	\end{itemize}
	\item If a lane cannot be used, due to poor quality or other reasons, skip it and proceed as described above.
\end{enumerate}

Example for the seeding of a 8-lane track \textit{from fastest to slowest seed time}:\\
100m and shorter: fastest time in lane 4, then lane 5,3,6,2,7,1,8\\
200m and longer: fastest time in lane 1, then lane 2,3,4,5,6,7,8

\section{Timing, Photo Finish and False Start Monitoring}

\begin{enumerate}[leftmargin=1.5em,labelsep=*]
	\item A Fully Automatic Timing and Photo Finish System must be used for the track races at Unicon and is 
  strongly recommended for track races at all other competitions.
	The system must have been tested, and have a certificate of accuracy issued within 4 years of the competition, 
  including the following:
	
	\begin{enumerate}[label*=\arabic*,leftmargin=1.5em,labelsep=*]
		\item The System must record the finish through a camera positioned in the extension of the finish line, 
    producing a composite photo finish image of at least 100 images per second, ideally 1000 images per second. 
    The image must be synchronized with a uniformly marked time-scale graduated in 0.01 seconds.
		
		\item The System must be started automatically by the Starter's signal, so that the overall 
    delay between the start signal and the start of the timing system is constant and equal to or 
    less than 0.001 second.
	\end{enumerate}
	\textbf{Note:} A system that works not automatically at start and finish will not produce fully 
  automatically measured times and therefore does not comply with this requirements.
	
	\item The placing and times of the riders must be read from the Photo Finish image by means 
  of a cursor with its reading line guaranteed to be perpendicular to the time scale.\\	
	\textbf{Note:} In order to confirm that the camera is correctly aligned and to facilitate the reading of 
  the Photo Finish image, the intersection of the lane lines and the finish line must be coloured black in a 
  suitable design. Any such design must be solely confined to the intersection, for no more than 20 mm beyond, 
  and not extended before, the leading edge of the finish line. Similar black marks may be placed on each side of 
  the intersection of an appropriate lane line and the finish line to further facilitate reading.
	
	\item The system must automatically determine and record the riders finish times and must be able to produce 
  a printed image (in physical form or into a file) showing the time of each rider. 
  Additionally, the system must provide a tabular overview showing the time or other result for each rider. 
  Subsequent changes of automatically determined values and manual input of values (like start time, finish time), 
  must be indicated by the system automatically in the time scale of the printed image and the tabular overview.
	
	\item For the track races at Unicon a false start monitoring system, which is able to reliably detect a 
  crossing of the start line before the start signal, must be used and is strongly recommended for track races 
  at all other competitions.
\end{enumerate}

\section{Accuracy of Results}

\begin{enumerate}[leftmargin=1.5em,labelsep=*]
	\item For all Gliding and Coasting disciplines where the distance is measured, unless the distance is an exact 
  0.1 meter, the distance must be converted and recorded to the next shorter 0.1 meter, e.g.\ 34.56 m 
  must be recorded as 34.5 m.
	When two riders reach the same distance, it must be determined to be a tie and the tie 
  must remain and gets published as such.
	
	\item For all Slow disciplines and Stillstand, unless the time is an exact 0.1 second, 
  the time must be converted and recorded to the next shorter 0.1 second.
	When two riders reach the same time, it must be determined to be a tie and the tie must 
  remain and gets published as such.
	
	\item For all other track racing events mentioned in this chapter, unless the time is an exact 0.01 second, 
  the time must be converted and recorded to the next longer 0.01 second, e.g.\ 14.533 seconds 
  must be recorded as 14.54 seconds.
	In the event that there is a tie where an award and/or a place in the final is at stake, 
  if the performances in question were achieved in the same heat and a photo finish system was used, 
  the image of this system must be used to decide on the placings.
	In this case, the note (Photo Finish: +0.00X) is printed on the results list next to the official time. 
  In other cases it must be determined to be a tie and the tie must remain and gets published as such.\\
  
	\textbf{Example:} If two riders have reached a time of 0:07.08 and the image of the photo finish 
  system shows a difference of 0.006 seconds, the following will be printed on the result list:\\
	\begin{tabular}{l l l}
	1st Place & Rider 1 & 0:07.08 \\
	2nd Place & Rider 2 & 0:07.08 (Photo finish: +0.006)\\
	\end{tabular}
\end{enumerate}

\section{Wind Measurement}

	It is recommended to measure the wind speed for 100 m, 200 m, One Foot, Wheel Walk, Track Coasting and 
  Track Gliding since tailwind can have a significant influence when the track is not 
  ridden at least once complete.
	For Unicons such wind speed measurement is required.
	For World Records to be valid, there may not be a tailwind averaging more than 2 m/s during the period 
  specified below.
	The following rules apply to wind measurement:

\begin{enumerate}[leftmargin=1.5em,labelsep=*]
	\item All wind gauge equipment must be IUF approved and manufactured and calibrated according to 
  international standards.
	The accuracy of the measuring eqipment used in the competition must have been verified by an 
  appropriate organisation accredited by the national measurement authority.\\
	\textbf{Note:} All WordAthletics certified devices are approved.
	
	\item Non-mechanical wind gauges must be used.
	
	\item The wind gauge should preferably be started and stopped automatically and remotely, 
  and the information conveyed directly to the competition computer.
	A manual start and stop should only be performed if there is no other possibility.
	In any case, the responsible juge needs a separate instruction for the correct operation of the wind gauge.
	
	\item The wind gauge must be read in meters per second, rounded to the next higher tenths of a meter 
  per second, unless the second decimal is zero, in the positive direction.
	Gauges that produce digital readings expressed in tenths of meters per second must be constructed so as 
  to comply with this rule.\\	
	\textbf{Explanation:} This means that a reading of +2.03 m/s must be recorded as +2.1 m/s; a reading of 
  -2.03 m/s must be recorded as -2.0 m/s.
	
	\item The wind gauge is placed beside the straight, adjacent to lane 1 and 50 m from the 
  finish line (200 m and 100 m) respectively 25 m from the finish line (One Foot and Wheel Walk).
	During the Track Coasting the wind gauge is placed beside the straight, adjacent to lane 1 and 50 m 
  from the starting line.
	In all cases the measuring plane must be positioned at 1.22 m ± 0.05 m height 
  and not more than 2 m away from the track.
	
	\item The period for which the wind velocity must be measured from the start signal are as follows:\\		
			100 m . . . . . . . . . . . . . . .  14 seconds,\\
			One Foot. . . . . . . . . . . . . .  8 seconds,\\
			Wheel Walk . . . . . . . . . . . . 8 seconds.\\
			For a 200 m race, the wind velocity must normally be measured for a period of 14 seconds commencing 
      when the first athlete enters the straight. The measuring can also be startet automatically after 
      12 seconds from the start signal.\\
			For Track Coasting the wind velocity must be measured for a period of 15 seconds from the time when 
      the athlete passes the start line.\\
\end{enumerate}